\scanpage{46}
\mspage{21}

Dans le cas où $X$ est compacte, on trouve que c'est la seule métrique riemannienne (comp. avec la structure conforme) qui soit de courbure constante égale à $-1$.

2) Cas \deleted{elliptique} comp. et (orientable ou non) de genre 1. Dans le cas orienté, i.e. celui d'une courbe elliptique complexe, dont le rev. universel est une droite complexe, on a sur celui-ci une métrique canonique mod scalaires (métrique hermitienne) invariante par translations (et qu. inv. par tout autom i.e. par transformation affine), qu'on peut descendre \unsure{à|?|!} $X$. La constante multiplicative est

% archive sheet number: 48
